% ============================================================
%  ch03 —— §2.1–2.3（前半，原文 p.12–p.15）
% ============================================================
\routenote{p.12--p.15}

\section{特征层与特征值——主要结果的陈述}

本节将构造自守层 $A_\varphi$。我们还将陈述关于 Kloosterman 层 $\Kl^{L}_G(\varphi,\chi)$ 的局部与整体单值性的主要结果；这个层将被定义为 $A_\varphi$ 的特征值。

\subsection{特征函数}

本小节给出 Gross 的自守表示 $\pi$ 中特征函数的一个简单公式。令人惊讶的是，这个计算不依赖于 Gross 的结果。我们将使用 §1.3 中选定的仿射一般特征 $\psi\circ\varphi:I(1)/I(2)\to\mathbb{Q}_\ell$。

由命题 1.1 知
\begin{gather*}
\Bun_{G(0,2)}(k)=G(k(t))\backslash G(0,2)(\mathbb{A}_{k(t)})\Big/\prod_x G(0,2)(\mathcal O_x)
=I^-(0)\backslash G(k((s)))/I(2),\\
G(k((s)))=\bigsqcup_{w\in\widetilde W}I^-(0)\,w\,I(1).
\end{gather*}
假设如引言中那样给定 $G(\mathbb{A}_{k(t)})$ 的一个自守表示 $\pi=\otimes'\pi_\nu$，使得对 $\nu\notin\{0,\infty\}$ 局部表示 $\pi_\nu$ 非分歧，$\pi_0$ 是 Steinberg 表示，且 $\pi_\infty$ 出现在 $\mathrm{c\text{-}Ind}_{Z(G(0,1))\times I(1)}^{G(k((s)))}\psi\circ\varphi$ 中。则存在 $\Bun_{G(0,2)}(k)$ 上的函数 $f$，使得对所有 $k\in I(1)$ 有 $f(gk)=\psi(\varphi(k))f(g)$。下面的引理以初等方式刻画这类函数：

\begin{lemma}
设 $f:\Bun_{G(0,2)}(k)\to\mathbb{Q}_\ell$ 是函数，使得对所有 $k\in I(1)$、$g\in G(0,2)(k((s)))$ 有 $f(gk)=\psi(\varphi(k))f(g)$。则 $f$ 由它在 $\gamma\in\Omega$ 处的值 $f(\gamma)$ 唯一决定。此外，对所有 $w\in\widetilde W-\Omega$ 有 $f(w)=0$。
\end{lemma}

\begin{proof}
对任意 $w\in\widetilde W$ 且 $l(w)>0$，存在单仿射根 $\alpha_i$ 使 $w.\alpha_i$ 为负。这蕴涵 $wU_{\alpha_i}w^{-1}\subset I^-(0)$。于是对所有 $u\in U_{\alpha_i}$ 有
\[
\psi(\varphi(u))f(w)=f(wu)=f(w),
\]
故 $f(w)=0$。对 $l(w)=0$ 的 $w$，即 $w\in\Omega$，由推论 1.3 可知 $f(\gamma)$ 的值可任意。
\end{proof}

\begin{remark}
上述引理中的任意函数 $f$ 自动是尖的。事实上，对任意抛物子群 $Q\subset G$ 及其幂幺根 $N_Q$，常数项函数
\[
f_{N_Q}(g)=\int_{N_Q(K)\backslash N_Q(\mathbb{A}_K)}f(ng)\,dn
\]
在 $N_Q(K_\infty)$ 下左不变，且相对于特征 $\psi\varphi$ 在 $I_\infty(1)$ 下右等变。与引理 2.1 类似的论证表明这样的函数必为零。
\end{remark}

\subsection{特征层}

让我们把前面的观察几何化。记 $\mathrm{AS}_\psi$ 为由特征 $\psi$ 定义的 $\Aff^1$ 上的 Artin--Schreier 层。令 $\mathrm{AS}_\varphi:=\varphi^*(\mathrm{AS}_\psi)$，即 Artin--Schreier 层到 $I(1)/I(2)\cong\Aff^d$ 的拉回。

记 $\mathrm{Perv}(\Bun_{G(0,2)})^{I(1),\mathrm{AS}_\varphi}$ 为 $\Bun_{G(0,2)}$ 上 $(I(1),\mathrm{AS}_\varphi)$-等变的 perverse 层范畴。对带 $I(1)/I(2)$ 作用的任意叠，我们也将一般地使用这个记号。

对任意 $\gamma\in\Omega$，记 $j_\gamma:I(1)/I(2)\hookrightarrow\Bun^\gamma_{G(0,2)}$ 为大胞腔的嵌入，$i_\gamma:\operatorname{Spec}(k)\to\Bun^\gamma_{G(0,2)}$ 为由 $\gamma$ 定义的 $G(0,2)$-丛给出的映射。下面是引理 2.1 的几何类比。

\begin{lemma}
层 $\mathrm{AS}_\varphi$ 满足 $j_{\gamma,!}\mathrm{AS}_\varphi=j_{\gamma,*}\mathrm{AS}_\varphi$。此外，对任意 $\gamma\in\Omega$，函子
\[
\mathrm{Perv}(\GG_m)\to\mathrm{Perv}\big(\Bun^\gamma_{G(0,2)}\times\GG_m\big),\qquad
F\mapsto j_{\gamma,!}\mathrm{AS}_\varphi\boxtimes F[\dim(\Bun_{G(0,2)})]
\]
是范畴等价。其逆由 $K\mapsto(i_\gamma\times\mathrm{id}_{\PP^1\setminus\{0,\infty\}})^*K[-\dim(\Bun_{G(0,2)})]$ 给出。
\end{lemma}

\begin{proof}
对任意 $w\in\widetilde W-\Omega$，在 $N(T((t)))$ 中取代表元，仍记为 $w$。考虑由 $w$ 定义的 $G(0,2)$-丛 $P_w$。设 $U_\alpha\subset I(1)$ 是对应于某个单仿射根 $\alpha$ 的根子群，使得 $w.\alpha$ 为负，即 $U_{w.\alpha}\subset I^-(0)$。这定义包含 $U_\alpha\hookrightarrow\operatorname{Aut}(P_w)$。于是得到交换图
\[
\begin{tikzcd}
U_\alpha\times\operatorname{Spec}(k) \arrow[r, "{(\mathrm{id},P_w)}"] \arrow[d] & \operatorname{Spec}(k) \arrow[d, "{P_w}"]\\
U_\alpha\times\Bun_{G(0,2)} \arrow[r, "{\text{作用}}"] & \Bun_{G(0,2)}.
\end{tikzcd}
\]
这蕴涵 $(\mathrm{AS}_\varphi|_{U_\alpha})\boxtimes K|_{P_w}\cong\mathbb{Q}_\ell|_{U_\alpha}\boxtimes K|_{P_w}$。由于我们假设 $\mathrm{AS}_\varphi|_{U_\alpha}$ 由 $U_\alpha$ 的非平凡特征定义，可知 $K$ 在 $P_w$ 处的茎为零。对偶地，同样的结果对 $K$ 在 $P_w$ 处的余茎成立。

这证明了第一个断言。对第二个断言，它还蕴涵：$\Bun^\gamma_{G(0,2)}$ 上任意 $(I(1),\mathrm{AS}_\varphi)$-等变 perverse 层都是其从子叠 $j_\gamma(I(1)/I(2))$ 的 $!$-扩张。在这个子叠上，与局部系统 $\mathrm{AS}_\varphi$ 作张量给出 $(I(1)/I(2))$-等变层与 $(I(1)/I(2),\mathrm{AS}_\varphi)$-等变层之间的等价。这证明了我们的断言。
\end{proof}

利用这个引理，我们现在可以定义自守层：

\begin{definition}
定义 $A_\varphi\in\mathrm{Perv}(\Bun_{G(0,2)})^{I(1),\mathrm{AS}_\varphi}$ 为在分支 $\Bun^\gamma_{G(0,2)}$ 上由 $j_{\gamma,!}\mathrm{AS}_\varphi[\dim(\Bun_{G(0,2)})]$ 给出的 perverse 层。记 $A^\gamma_\varphi$ 为 $A_\varphi$ 到分支 $\Bun^\gamma_{G(0,2)}$ 的限制。
\end{definition}

\begin{remark}[带乘性特征的变体]
我们可以把 $A_\varphi$ 的上述构造稍作推广。回忆 $\Bun^\gamma_{G(1,2)}$ 中的开胞腔典范同构于 $T_0^{\mathrm{red}}\times I(1)/I(2)$\zhnote{原文此处写作 $I(1)/I(0)$；由推论 1.3(6) 与 $I(1)/I(2)\cong\Aff^n$ 的记号，应为 $I(1)/I(2)$，属原文笔误。以下同。[来源：原文 p.13 与推论 1.3(6)]}。任意特征 $\chi:T_0^{\mathrm{red}}(k)\to\mathbb{Q}_\ell$ 在 $T_0^{\mathrm{red}}$ 上定义特征层 $\mathrm{Kum}_\chi$。于是上述引理 2.3 对 $\Bun_{G(1,2)}$ 上 $(T_0^{\mathrm{red}}\times I(1)/I(2),\mathrm{Kum}_\chi\boxtimes\mathrm{AS}_\varphi)$-等变层也成立。记 $\Bun_{G(1,2)}$ 上相应的 perverse 层为 $A_{\varphi,\chi}$。
\end{remark}

\subsection{几何 Hecke 算子}

为陈述主要结果，我们需要回忆几何 Hecke 算子的定义。

Hecke 修改之叠是
\[
\begin{aligned}
\Hecke_{\mathbb{G}(m,n)}^{\Aff^1}(S):=\Big\{(E_1,E_2,x,\phi)\ \Big|\ &E_i\in\Bun_{\mathbb{G}(m,n)}(S),\ x:S\to\Aff^1,\\
&\phi:E_1|_{\PP^1-x\times S}\xrightarrow{\sim}E_2|_{\PP^1-x\times S}\Big\}.
\end{aligned}
\]
这个叠有自然的遗忘映射：
\begin{equation}
\begin{tikzcd}[column sep=large]
& \Hecke^{\Aff^1}_{\mathbb{G}(m,n)} \arrow[dl, "\mathrm{pr}_1"'] \arrow[dr, "\mathrm{pr}_2"] \arrow[d, "\mathrm{pr}_{\Aff^1}"] & \\
\Bun_{\mathbb{G}(m,n)} & \Bun_{\mathbb{G}(m,n)}\times\Aff^1 \arrow[l, "\mathrm{proj}_1"'] \arrow[r] & \Aff^1
\end{tikzcd}
\end{equation}
其中 $\mathrm{pr}_1(E_1,E_2,x,\phi)=E_1$，$\mathrm{pr}_2(E_1,E_2,x,\phi)=(E_2,x)$，$\mathrm{pr}_{\Aff^1}(E_1,E_2,x,\phi)=x$。记
\[
\Hecke_{\mathbb{G}(m,n)}^{\PP^1\setminus\{0,\infty\}}:=\mathrm{pr}_{\Aff^1}^{-1}\big(\PP^1\setminus\{0,\infty\}\big).
\]

\begin{remark}\label{rem:2.7}
\begin{enumerate}[label=(\arabic*), leftmargin=1.8em, itemsep=.2em]
\item $\mathrm{pr}_1$ 在平凡丛 $\mathbb{G}(m,n)\in\Bun_G$ 上的纤维称为 Beilinson--Drinfeld Grassmannian，记作 $\mathrm{GR}_{\mathbb{G}(m,n)}$。$\mathrm{GR}_{\mathbb{G}}\to\Aff^1$ 在点 $x\in\PP^1\setminus\{0,\infty\}$ 上的纤维同构于仿射 Grassmannian $\mathrm{Gr}_{\mathbb{G},x}$，即商 $\mathbb{G}(K_x)/\mathbb{G}(\mathcal O_x)$。
\item $\mathrm{pr}_2$ 在 $\Bun_{\mathbb{G}(m,n)}\times\PP^1\setminus\{0,\infty\}$ 上的几何纤维（非典范地）同构于仿射 Grassmannian $\mathrm{Gr}_{\mathbb{G}}$。在 $\Bun_{\mathbb{G}}\times\PP^1\setminus\{0,\infty\}$ 的光滑拓扑下局部，这个纤维化是平凡的（例如注 4.1）。
\item 图 (2.6) 有一大族对称性。群 $I(0)/I(2)=(I(1)/I(2))\rtimes T^{\mathrm{red}}_\infty$ 通过改变 $\infty$ 处的 $I(2)$-水平结构作用在 $\Bun_{\mathbb{G}(m,2)}$ 上，且这个作用扩充到图 (2.6)（即它也作用在 $\Hecke^{\Aff^1}_{\mathbb{G}(m,2)}$ 上，且映射 $\mathrm{pr}_i$ 在这些作用下等变）。一维环面 $\GG_m^{\mathrm{rot}}$ 作用在固定点 $\{0\}$ 与 $\{\infty\}$ 的曲线 $\PP^1$ 上，故也作用在 (2.6) 上。最后，钉住自同构 $\operatorname{Aut}^\dagger(\mathbb{G})$ 作用在 (2.6) 上。于是我们看到群
\[
I(0)/I(2)\rtimes\big(\GG_m^{\mathrm{rot}}\times\operatorname{Aut}^\dagger(\mathbb{G})\big)
\]
作用在图 (2.6) 上。
\end{enumerate}
\end{remark}

先回忆常群概形的 Hecke 算子。为此我们收集几何 Satake 等价的一些事实（见 [30]、[18]）。

设 $\mathrm{Gr}_G=G((\tau))/G[[\tau]]$ 为抽象的仿射 Grassmannian，不参照 $\PP^1\setminus\{0,\infty\}$ 上的任何点。令 $\mathcal O=k[[\tau]]$，$\operatorname{Aut}_{\mathcal O}$ 为 $\mathcal O$ 的连续（在 $\tau$-进拓扑下）自同构的 pro-代数 $k$-群。则 $G[[\tau]]\rtimes\operatorname{Aut}_{\mathcal O}$ 从左侧作用在 $\mathrm{Gr}_G$ 上。$G[[\tau]]$ 在 $\mathrm{Gr}_G$ 上的轨道由支配余特征 $\mu\in X_*(T)^+$ 编号。轨道记为 $\mathrm{Gr}_{G,\mu}$，其闭包（Schubert 簇）记为 $\mathrm{Gr}_{G,\le\mu}$。记 $\mathrm{Gr}_{G,\le\mu}$ 的相交上同调层为 $\mathrm{IC}_\mu$。我们将把 $\mathrm{IC}_\mu$ 归一化为权 $0$（见下面的注 2.9）。

Satake 范畴 $\mathrm{Sat}=\mathrm{Perv}_{\operatorname{Aut}_{\mathcal O}}(G[[\tau]]\backslash G((\tau))/G[[\tau]])$ 是 $\mathrm{Gr}_G$ 上 $G[[\tau]]\rtimes\operatorname{Aut}_{\mathcal O}$-等变 perverse 层（带有限型支撑）的范畴。类似地，把情形基变换到 $\bar k$ 定义 $\mathrm{Sat}^{\mathrm{geo}}$。最后把规范化的半单 Satake 范畴 $\mathcal S$ 定义为 $\mathrm{Sat}$ 中由 $\mathrm{IC}_\mu$ 的直和组成的满子范畴。

[30] 与 [18] 证明了：
\begin{itemize}[leftmargin=1.6em, itemsep=.2em]
\item $\mathrm{Sat}^{\mathrm{geo}}$ 带有自然的张量结构（它对 $\mathrm{Sat}$ 也有定义），使得整体上同调函子 $h=H^*(\mathrm{Gr}_G,-):\mathrm{Sat}^{\mathrm{geo}}\to\mathrm{Vec}$ 是纤维函子。
\item $\operatorname{Aut}^\otimes(h)$ 是 $\mathbb{Q}_\ell$ 上连通约化群，它 Langlands 对偶于 $G$。令 $\ckG:=\operatorname{Aut}^\otimes(h)$，则 Tannaka 形式化给出张量范畴的几何 Satake 等价
\[
\mathrm{Sat}^{\mathrm{geo}}\cong\Rep(\ckG).
\]
\item 按构造，$\ckG$ 配有极大环面 $\ckT$ 与自然同构 $X^*(\ckT)\cong X_*(T)$（事实上 $\ckG$ 配有典范钉住；见引理 9.3）。
\end{itemize}

几何 Satake 等价把 $\mathrm{IC}_\mu$ 送到极值权为 $\mu$ 的不可约表示 $V_\mu$。记这个等价之逆为 $V\mapsto\mathrm{IC}_V$。

\begin{remark}
在 [1, §3.5] 中论证了 $\mathcal S$ 在 $\mathrm{Sat}$ 的张量结构下封闭。因此 $\mathcal S$ 自然是张量范畴。

沿 $\mathrm{Gr}_G\otimes_k\bar k\to\mathrm{Gr}_G$ 的拉回给出张量函子 $\mathcal S\to\mathrm{Sat}^{\mathrm{geo}}$；由于两个范畴都是半单的且单对象显式已知，容易看出它是等价。因此 [30] 与 [18] 中的上述结果都适用于 $\mathcal S$。特别地，我们得到（$k$-版本的半单化）几何 Satake 等价
\[
\mathcal S\cong\mathrm{Sat}^{\mathrm{geo}}\cong\Rep(\ckG).
\]
\end{remark}

\begin{remark}[权的归一化]
我们采用使该复形纯权 $0$ 的归一化，即在 $\mathbb{Q}_\ell$ 中取 $q$ 的一个平方根，并把 $\mathrm{IC}_\mu$ 记为相交复形张量上 $\mathbb{Q}_\ell\big(\tfrac12\dim(\mathrm{Gr}_\mu)\big)$。

正如 [15] 所指出的，并不必要作这个相当不自然的选择——作它是为了从范畴 $\mathcal S$ 得到群 $\ckG$。替代地，我们可以把范畴 $\mathcal S$ 扩大，纳入相交上同调层 $\mathrm{IC}_\mu(n)$ 的所有 Tate 扭转。由前一注，这仍是一个中性张量范畴，定义群 $\ckG_1$，它是 $\ckG$ 被一个中心一维环面所作的扩张。
\end{remark}
